
%------------------------------------Outer Title Page--------------------------------
\begin{titlepage}
\begin{center}\sf \Large
{\Huge\textbf{\thesistitle}}

\vfill
{\baselineskip=25pt
A dissertation submitted by\\

{\bf  \authorname}\\

(Reg. No. 201721410)

\vspace{20pt}
in partial fulfillment\\
for the award of the Degree of \\

{\bf Master of Science}\\

(Mathematics)
\vfill

% \vspace{50pt}
% \hspace{-2in}
% \begin{center}
\includegraphics[height=4cm]{smit.eps}

% \flushright
% \begin{minipage}{60pt}
% \includegraphics[height=1.5cm]{manipal.eps}
% \end{minipage}

% \end{center}



DEPARTMENT OF MATHEMATICS\\
SIKKIM MANIPAL INSTITUTE OF TECHNOLOGY\\
\today}
\end{center}
\end{titlepage}
%------------------------------------End of the Outer Title Page--------------------------------

\begin{preface}{\sf Declaration}
% \thispagestyle{plain}{
% \fancyhf{}
% \fancyfoot[C]{}
% \fancyfoot[C]{{\sf{--~\thepage~--}}} 
% \renewcommand{\headrulewidth}{0.0pt}}
\baselineskip=18pt
\vspace*{0.8in}
I do hereby declare that the work contained in this dissertation entitled ``{\bf \thesistitle}'' 
has done by me, under the supervision of Dr. Biswajit Deb,
Assistant Professor, Department of Mathema\thispagestyle{plain}{
\fancyhf{}
\fancyfoot[C]{}
\fancyfoot[C]{{\sf{--~\thepage~--}}} 
\renewcommand{\headrulewidth}{0.0pt}}tics, SMIT in partial fulfillment
for the award of the degree of Master of Science and
this work has not been submitted elsewhere for a degree or diploma.

\vspace*{1.5in}
\begin{minipage}[t]{0.3\textwidth}
\hspace*{-18pt}\today
\end{minipage}
\begin{minipage}[t]{0.65\textwidth} \thispagestyle{plain}{
% \fancyhf{}
% \fancyfoot[C]{}
% \fancyfoot[C]{{\sf{--~\thepage~--}}} 
% \renewcommand{\headrulewidth}{0.0pt}}
\begin{flushright}
\baselineskip=18pt
    {\bf    \authorname}\\
    Reg No.201721410\\
    Department of Mathematics\\
    Sikkim Manipal Institute of Technology\\
    India.
\end{flushright}
\end{minipage}
\end{preface}

%%%%%%%%%%%%%%%%%%%%%%%%%%%%%%%%%%%%%%%%%%%%%%%%%%%%%%%%%%%%%%%%%%%%%%%%%%%%

\begin{preface}{\sf Certificate}
% \thispagestyle{plain}{
% \fancyhf{}•••••••••••
% \fancyfoot[C]{}
% \fancyfoot[C]{{\sf{--~\thepage~--}}} 
% \renewcommand{\headrulewidth}{0.0pt}}
\baselineskip=18pt
\vspace*{0.8in}
This is to certify that the  dissertation entitled
``{\bf \thesistitle}'' by  \authorname,  a student of Department of
Mathematics, SMIT, Sikkim in partial fulfillment for the award of the degree
of Master of Science has been carried out under the supervision of Dr. Biswajit Deb and that this
work has not been submitted elsewhere for a degree or diploma.


\vspace*{1.5in}
\noindent

\begin{minipage}[t]{0.50\textwidth}
\baselineskip=18pt
    \noindent {\bf Prof. (Dr.) Anjan Raychaudhuri}\\
    Head of the Department\\
    Department of Mathematics\\
    SMIT, India
    \end{minipage}
    \hspace{.4in}
    \begin{minipage}[t]{0.65\textwidth}
%     \begin{flushright}
\baselineskip=18pt
% \hspace{.2in}{
  {\bf Biswajit Deb}\\
    Assciate Professor\\
    Department of Mathematics\\
    SMIT, India
% \end{flushright}
\end{minipage}
\end{preface} % Certificate
%%%%%%%%%%%%%%%%%%%%%%%%%%%%%%%%%%%%%%%%%%%%%%%%%%%%%%%%%%%%%%%%%%%%%%%%%%%%%%%%%%%%%%%%%%%%%%%%%%%%
\begin{preface}{\sf Acknowledgements}
% \thispagestyle{plain}{
% \fancyhf{}
% \fancyfoot[C]{}
% \fancyfoot[C]{{\sf{--~\thepage~--}}} 
% \renewcommand{\headrulewidth}{0.0pt}}
\baselineskip=18pt
 
Foremost, I would like to express my special thanks of gratitude to my guide Dr.Biswajit
Deb, Department of Mathematics, SMIT, for the valuable guidance and constant support during this journey. I could not have 
imagined having better guide and mentor for my Project.

Beside my guide, I would also like to express my special thanks to Prof. (Dr.) Anjan Raychaudhuri, Head of the Department, for 
encouragement and valuable suggestion.
 
My special thanks to all the faculty members of the Department of
Mathematics, SMIT, for their overall support.
I thank my fellow mates for the suggestion they gave and for all the fun we had during last two years.


Last but not the least, I would like to thank my mother and sisters for supporting me spiritually and mentally throughout my life.


\vspace{2.5cm}
\noindent
\begin{minipage}[t]{0.3\textwidth}
\hspace*{1pt}\today
\end{minipage}
\begin{minipage}[t]{0.65\textwidth}
\begin{flushright}
\hspace{0pt}\baselineskip=18pt
    {\bf  \authorname}\\
    SMIT, Sikkim\\
\end{flushright}
\end{minipage}
%%%%%%%%%%%%%%%%%%%%%%%%%%%%%%%%%%%%%%%%%%%%%%%%%%%%%%%%%%%%%%%%%%%%%%%%%%%%%%%%%%%%%%%%%%%%%%%%%
\end{preface} % Acknowledgements
\begin{preface}{\sf Abstract}
% \thispagestyle{plain}{
% \fancyhf{}
% \fancyfoot[C]{}
% \fancyfoot[C]{{\sf{--~I~--}}} 
% \renewcommand{\headrulewidth}{0.0pt}}
\baselineskip=18pt•••••••••••
\addcontentsline{toc}{chapter}{Abstract}

A marked graph is a graph with a $+$ or $-$ sign on each vertex and is called consistent if each cycle has an even number of $−$ 
signs. This concept is motivated by problems of communication networks and social networks. Formally, a marked graph is a graph 
$G(V,E)$ together with a function $\mu~:~V\rightarrow\{-1,1\}$. The function $\mu$ is known as the 
marking of the marked graph and for any $v\in V$, we call $\mu(v)$ to be the mark of $v$. By mark of a cycle in a marked graph we 
mean the product of the mark of the vertices in the cycle. A marked graph in which every cycle has a positive mark is called a 
consistent marked graph.
A signed graph $\Sigma=(G,\sigma)$ is a graph $G(V,E)$ together with a function $\sigma 
~:~E\rightarrow\{-1,1\}$. For any signed graph $\Sigma=(G,\sigma)$,
the marking $\mu~:~V\rightarrow\{-1,1\}$ defined by 
$$\mu (v)=\begin{cases}
+1, & \text{ if $v$ is isolated;} \\
\prod\limits_{u\in N(v)}\sigma(uv), & \text{ otherwise;} 
\end{cases}
$$
  where $N(v)$ is the open neighborhood of $v$ in $G$ is said to be the canonical marking.  In this 
dissertation we have  investigated the relationship of balancedness and canonical consistency of a signed 
graph to its 2-path signed graphs.

\vspace{1.5cm}

\end{preface} % abstract
